\documentclass[10pt,a4paper]{amsart}
\usepackage[margin=19mm]{geometry}
\usepackage{amsmath,amssymb,mathtools}
\usepackage[scr=rsfs,cal=euler]{mathalfa}
\usepackage{xfrac}
\usepackage[unicode,hidelinks,bookmarks=false]{hyperref}
\newcommand{\ie}{i.e.\ }
\newcommand{\cf}{cf.\ }
\newcommand{\resp}{respectively\ }
\newcommand{\Subsection}{\subsection}
\providecommand{\nobreakdash}{\nobreak\hbox{-}}
\setlength{\emergencystretch}{2em}
\multlinegap0pt
\usepackage{bm}
\usepackage{bbm}
\usepackage{dsfont}
\usepackage{xspace}
\usepackage{cancel}
\usepackage{mathtools}
\usepackage{amsthm}
\makeatletter
\@ifundefined{defn}{\newtheorem{defn}{Defn}}{}
\@ifundefined{thm}{\newtheorem{thm}[defn]{Theorem}}{}
\makeatother
\title[Mean Li-Yorke chaos for a sequence of operators on Banach sp]{Mean Li-Yorke chaos for a sequence of operators on Banach spaces}
\author{Jian Li, Xinsheng Wang, Jianjie Zhao}
\date{Source: arXiv:2511.06207v1 (CC BY 4.0)}
\begin{document}
\maketitle

\noindent\emph{Source and licence.} Mean Li-Yorke chaos for a sequence of operators on Banach spaces. Version arXiv:2511.06207v1 (CC BY 4.0), distributed under the Creative Commons Attribution 4.0 International licence, as recorded in the arXiv licence metadata for this exact version and shown on the abstract page https://arxiv.org/abs/2511.06207v1. Full rights notice: licenses/arxiv-2511.06207-CC-BY-4.0.txt.\par\smallskip

\noindent\textit{Editorial scope.} Counted unit: the Thm whose source label is \texttt{thm:mean-sensitive} in the article (source0.tex lines 255--291, source label \texttt{thm:mean-sensitive}), delivered together with the article's own complete original proof of it (source0.tex lines 293--408, one \texttt{proof} environment of 116 lines). No mathematics is written, repaired or re-derived: statement and proof are the article's own text line for line. Required context retained uncounted: thm:mean-UBT (lines 168--188). Every definition, display and label of those lines is retained unchanged. Omitted from this package and not redistributed: 1--167, 189--254, 292--292, 409--1538. Nothing inside a retained excerpt is omitted or altered: every retained body is delivered exactly as the article's lines are written, so no omission receipt of any kind accompanies this package. Citations: the retained excerpts carry no citation command at all, so no bibliography entry of the article is transcribed and no reference is redistributed. Reference adaptation: none was needed and none was made; every reference inside the delivered unit resolves inside it, so no unresolved reference or citation is produced. Printed slips kept literal: no printed word, symbol, hypothesis or number of the counted statement or of its proof was altered. Layout and class: the article's own class is \texttt{amsart} with packages including fontenc, geometry, amssymb, mathtools, newpxtext, newpxmath, fancyhdr, hyperref, amsrefs; the wrapper is the lane's pinned \texttt{amsart} editorial wrapper at 10pt on A4, which restates the theorem-like environments the retained text needs and adds the article's own macro definitions that the retained text needs (none), with extra wrapper packages (bm, bbm, dsfont, xspace, cancel, mathtools). The delivered numbering is the unit's own. Assets: no retained span contains a figure environment, a graphics-inclusion command, a PDF-page inclusion, a table or a TikZ picture; no image file is copied into this package, the package declares no asset dependency and no image file is redistributed with it. Licence: version arXiv:2511.06207v1 of this article is distributed under the Creative Commons Attribution 4.0 International licence, as recorded in the arXiv licence metadata for this exact version and shown on the abstract page https://arxiv.org/abs/2511.06207v1. A full rights notice is delivered at licenses/arxiv-2511.06207-CC-BY-4.0.txt. Characters: every retained excerpt is pure ASCII, so the delivered engine, XeLaTeX with its default Latin Modern text font, reports no missing character.
\begin{thm}\label{thm:mean-UBT}
    Let $(T_i)_{i=1}^{\infty}$ be a sequence of bounded linear operators from a Banach space $X$ to a normed linear space $Y$.
    Then the following assertions are equivalent:
    \begin{enumerate}
        \item \label{thm1: ME}
              the sequence $(T_i)_{i=1}^{\infty}$  is mean equicontinuous;
        \item \label{thm2: x bounded}
             for all $x\in X$,
              \[
                  \sup_{n\in\mathbb{N}}\frac{1}{n}\sum_{i=1}^{n}\Vert T_i x\Vert<\infty;
              \]
        \item \label{thm3: aCb}
             the sequence $(T_i)_{i=1}^{\infty}$ is absolutely Ces\`aro bounded;
        \item \label{thm4: ME-sup}
             for every $\varepsilon>0$, there exists $\delta>0$ such that for every $x,y\in X$ with $\Vert x-y\Vert<\delta$,
              \[
                  \sup_{n\in\mathbb{N}}\frac{1}{n}\sum_{i=1}^{n}\Vert T_ix-T_iy\Vert<\varepsilon.
              \]

    \end{enumerate}
\end{thm}

\begin{thm} \label{thm:mean-sensitive}
    Let $(T_i)_{i=1}^{\infty}$ be a sequence of bounded linear operators from a Banach space $X$ to a normed linear space $Y$.
    Then the following assertions are equivalent:
    \begin{enumerate}
        \item \label{thm1: mean sensitive}
              the sequence $(T_i)_{i=1}^{\infty}$  is mean sensitive;
        \item \label{thm2: x unbounded}
              there exists $x\in X$ such that
              \[
                  \limsup_{n\to{\infty}}\frac{1}{n}\sum_{i=1}^{n}\Vert T_i x\Vert=\infty;
              \]
        \item \label{thm3: residual points in X}
              the collection
              \[
                  \biggl \{ x\in X\colon \limsup_{n\to{\infty}}\frac{1}{n}\sum_{i=1}^{n} \Vert T_i x\Vert=\infty \biggr \}
              \]
              is residual in $X$;

        \item \label{thm4: residual points in X*X}
              the collection
              \[
                  \biggl\{ (x,y)\in X\times X\colon \limsup_{n\to{\infty}}\frac{1}{n}\sum_{i=1}^{n}\Vert T_i (x-y)\Vert=\infty \biggr \}
              \]
              is residual in $X\times X$;

        \item \label{thm5: bounded sequence}
              there exists a bounded sequence  $(y_k)_k$ in $X$ and a sequence $(N_k)_k$ in $\mathbb{N}$ such that
              \[
                  \sup_k\frac{1}{N_k}\sum_{i=1}^{N_k}\Vert T_i y_k\Vert=\infty;
              \]
        \item \label{thm6: sequence to 0}
              there exists a sequence  $(y_k)_k$ in $X$ and a sequence $(N_k)_k$ in $\mathbb{N}$ such that $\lim_{k\to\infty}y_k=0$ and
              \[
                  \liminf_{k\to\infty} \frac{1}{N_k}\sum_{i=1}^{N_k}\Vert T_i y_k\Vert>0.
              \]
    \end{enumerate}
\end{thm}

\begin{proof}
     $(\ref{thm1: mean sensitive}) \Rightarrow (\ref{thm2: x unbounded})$. 
    Otherwise, if for any $x\in X$,
    $$\limsup_{n\to \infty} \frac{1}{n} \sum_{i=1}^{n} \Vert T_ix\Vert<\infty,$$
    by Theorem \ref{thm:mean-UBT}, $(T_i)_{i=1}^{\infty}$ is mean equicontinuous. 
    Let $\delta$ be the constant in the definition of mean sensitivity.
    On the one hand, mean equicontinuity implies that there exists $\delta'>0$
    such that for every $x,y\in X$ with $\Vert x-y\Vert<\delta'$, one has
    \[
        \limsup_{n\to \infty} \frac{1}{n} \sum_{i=1}^{n} \Vert T_ix -T_iy\Vert <\delta.
    \]
    On the other hand, mean sensitivity implies that for every $x\in X$ and above $\delta'$,
    there exists some $y\in X$ with $\Vert x-y\Vert <\delta'$ such that
    \[
        \limsup_{n\to \infty} \frac{1}{n} \sum_{i=1}^{n} \Vert T_ix -T_iy\Vert >\delta.
    \]
    Which is a contradiction.
    
     $(\ref{thm2: x unbounded})  \Rightarrow (\ref{thm3: residual points in X})$. 
    Assume that $x_0\in X$ such that
    $$\limsup_{n\to \infty} \frac{1}{n}
        \sum_{i=1}^{n} \Vert T_ix_0\Vert=\infty.$$
    Let \[
        Y=\biggl\{ x\in X\colon \limsup_{n\to{\infty}}\frac{1}{n}\sum_{i=1}^{n}\Vert T_i x\Vert=\infty \biggr \}.
    \]
    It is clear that
    \[
        Y =\bigcap_{k=1}^\infty \biggl\{ x\in X\colon \exists n\geq k \text{ s.t. }
        \frac{1}{n}\sum_{i=1}^{n}\Vert T_i x\Vert >k\biggr\}.\]
    Then $Y$ is a $G_\delta$ subset of $X$.
    It suffices to show that for any $x\in X$ and $\varepsilon>0$, there exists some $y\in Y$ such that
    $\Vert x-y \Vert <\varepsilon$.
     
    If $x\in Y$, there's nothing to prove.
    If $x\notin Y$, put $y=x+\frac{\varepsilon}{2\Vert x_0\Vert}x_0 $,
    then $\Vert x-y \Vert <\varepsilon$
    and
    \begin{align*}
        \limsup_{n\to \infty} \frac{1}{n} \sum_{i=1}^n
        \Vert T_iy\Vert & =
        \limsup_{n\to \infty} \frac{1}{n} \sum_{i=1}^n \Vert T_i(x+\tfrac{x_0}{2\Vert x_0\Vert } \varepsilon )\Vert  \\
        & \ge \limsup_{n\to \infty} \frac{1}{n} \sum_{i=1}^n \Vert T_i( \tfrac{x_0}{2\Vert x_0\Vert } \varepsilon )\Vert -\limsup_{n\to \infty} \frac{1}{n} \sum_{i=1}^n \Vert T_i x \Vert \\
        & =\frac{ \varepsilon}{2\Vert x_0\Vert }  \limsup_{n\to \infty} \frac{1}{n} \sum_{i=1}^n \Vert T_i x_0  \Vert -\limsup_{n\to \infty} \frac{1}{n} \sum_{i=1}^n \Vert T_i x \Vert    \\
         & =\infty.
    \end{align*}

     $(\ref{thm3: residual points in X}) \Rightarrow (\ref{thm4: residual points in X*X})$.  It follows from the linearity of $(T_i)_i$.

     $ (\ref{thm4: residual points in X*X}) \Rightarrow (\ref{thm5: bounded sequence})$.  It is obvious.

     $(\ref{thm5: bounded sequence}) \Rightarrow (\ref{thm6: sequence to 0})$. 
    Assume that $(z_k)_k$ is the bounded sequence in $X$ and
    $(N_k)_k$ is the sequence in $\mathbb{N}$ such that
    \[
        \sup_k \frac{1}{N_k}\sum_{i=1}^{N_k} \Vert T_iz_k \Vert =\infty.
    \]
    Let $y_k=\frac{N_k}{ \sum_{i=1}^{N_k} \Vert T_iz_k \Vert} z_k$,
    then $\lim_{k\to \infty}y_k=0$ and
    \[
        \liminf_{k\to \infty} \frac{1}{N_k}\sum_{i=1}^{N_k} \Vert T_iy_k \Vert =1>0.
    \]

     $ (\ref{thm6: sequence to 0}) \Rightarrow (\ref{thm1: mean sensitive})$. 
    Suppose the sequences $(y_k)_k$ and $(N_k)_k$ are given by condition $(\ref{thm6: sequence to 0})$, and let
    \[
        \delta=\liminf_{k\to\infty} \frac{1}{N_k}\sum_{i=1}^{N_k}\Vert T_i y_k\Vert.
    \]
    Without loss of generality, we may assume that $(N_k)_k$ is increasing.
    For each $n\in \mathbb{N}$, let
    \[
        X_n=\biggl\{ x\in X: \exists k>n \ \text{s.t.} \ \frac{1}{k}\sum_{i=1}^{k} \Vert T_ix \Vert >
        \frac{\delta}{2} -\frac{1}{n} \biggr\}.
    \]
    Then each $X_n$ is an open subset of $X$.
    We will show that it is also dense in $X$.
    Let $U$ be a nonempty open subset of $X$ and pick $x\in U$.
    If $x\notin X_n$, then for any $k>n$, one has
    \[
        \frac{1}{k}\sum_{i=1}^{k}\Vert T_ix \Vert \le \frac{\delta}{2} -\frac{1}{n}.
    \]
    Since
    \[
        \liminf_{k\to \infty} \frac{1}{N_k} \sum_{i=1}^{N_k} \Vert T_iy_k \Vert =\delta,
    \]
    there exists $N\in \mathbb{N}$ such that for all $k>N$,
    we have
    \[
        \frac{1}{N_k} \sum_{i=1}^{N_k} \Vert T_iy_k \Vert > \delta -\frac{1}{2n}.
    \]
    As $\lim_{k\to \infty} (y_k+x) = x\in U$,
    there exists $k_0\in \mathbb{N}$ with $k_0>N$ such that $y_{k_0}+x \in U$
    and $N_{k_0}>n$.
    Then
    \begin{align*}
        \frac{1}{N_{k_0}}\sum_{i=1}^{N_{k_0}}\Vert T_i (y_{k_0}+x)\Vert
         & \ge \frac{1}{N_{k_0}}\sum_{i=1}^{N_{k_0}}\Vert T_i y_{k_0}\Vert  -\frac{1}{N_{k_0}}\sum_{i=1}^{N_{k_0}}\Vert T_i x\Vert \\
         & \ge \delta-\frac{1}{2n} - \Bigl(\frac{\delta}{2} - \frac{1}{n}\Bigr)                                                    \\
         & =\frac{\delta}{2} + \frac{1}{2n} > \frac{\delta}{2}-\frac{1}{n}.
    \end{align*}
    Which implies that $y_{k_0}+x\in U\cap X_n$.
    Therefore, the set
    \[
        X_0= \biggl\{x\in X: \limsup_{n\to \infty} \frac{1}{n}\sum_{i=1}^{n}\Vert T_ix \Vert \ge \frac{\delta}{2} \biggr\} =\bigcap_{n=1}^{\infty}X_n
    \]
    is residual in $X$.
    Fix $x\in X$ and $\varepsilon>0$,
    since $X_0$ is residual, there exists $z\in X$ with $\Vert z \Vert <\varepsilon$ and
    \[
        \limsup_{n\to \infty} \frac{1}{n}\sum_{i=1}^{n}\Vert T_iz \Vert \ge \frac{\delta}{2}.
    \]
    Let $y=z+x$, then $\Vert x-y \Vert <\varepsilon$ and
    \[
        \limsup_{n\to \infty} \frac{1}{n}\sum_{i=1}^{n}\Vert T_ix -T_iy\Vert =\limsup_{n\to \infty} \frac{1}{n}\sum_{i=1}^{n}\Vert T_iz \Vert \ge \frac{\delta}{2}.
    \]
    This shows that the sequence $(T_i)_{i=1}^{\infty}$ is mean sensitive.
\end{proof}

\end{document}
